%% ============================================================
%% sections/s2-findings.tex
%% H. R. 9510 Bill v5.0 - SEC. 2. FINDINGS; THE EVIDENCE TO LAW TO COST RECORD.
%% Rendered visuals: Figure 1 (gray-scale Mermaid lineage, TikZ), Table 1 (the
%% four works with the financial footprint), Figure 2 (ASCII cost ledger). The
%% carried v3.0 findings are condensed into a lettered narrative and extended
%% by the four v5.0 financial findings.
%% ============================================================

\billsec{SEC. 2.}{FINDINGS; THE EVIDENCE TO LAW TO COST RECORD.}\phantomsection\label{toc:s2}

\uscprov{1}{\textbf{(a) In General.} Congress finds that an automated process
must verify robot-patient interaction code before that code is generated or
executed in a Physical AI oncology clinical investigation; that the record of
the VVUQ-01 through VVUQ-04 developments establishes both the need for and the
feasibility of that requirement; and that the same record, read as a ledger,
establishes a third thing the prior versions left implicit: the requirement is
affordable, and the verification data that proves safety becomes financial
data the moment a cost is attached to each run \cite{kawchak2026national,
kawchak2026vvuq01, kawchak2026vvuq02}. Figure 1 shows how engineering evidence
became proposed law, then a precise amendment, and now a priced one.}

\begin{mermaidfig}
\node[mmin]   (v1) at (0,0)      {VVUQ-01\\method + 51-test pipeline};
\node[mmin]   (v2) at (4.0,0)    {VVUQ-02\\10-gate humanoid, 172 tests};
\node[mmstep] (v3) at (8.0,0)    {VVUQ-03\\standalone bill (v1.0)};
\node[mmstep] (v4) at (0,-2.4)   {VVUQ-04\\FD\&C amendment (v2.0)};
\node[mmlaw]  (v5) at (4.0,-2.4) {v3.0 visual and\\v4.0 Mermaid amendments};
\node[mmgoal] (b)  at (8.0,-2.4) {H. R. 9510 v5.0\\the Financial Data\\Amendment};
\draw[mmedge] (v1) -- (v2);
\draw[mmedge] (v2) -- (v3);
\draw[mmedge] (v3.south) -- ++(0,-0.55) -- ++(-8.0,0) -- (v4.north);
\draw[mmedge] (v4) -- (v5);
\draw[mmedge] (v5) -- (b);
\end{mermaidfig}
\vspace{-0.7cm}
\figcaption{Figure 1. Evidence-to-law-to-cost lineage: the method and pipeline, the hard\\engineering proof, the bill, the
amendment, the two visual versions, and this financial\\version, which attaches
a price to everything the prior five versions recorded.}

\uscprov{1}{\textbf{(b) The Embodiment Problem and the Sequencing Gap.}
Physical AI oncology investigations place autonomous and semi-autonomous
robots in direct physical contact with patients, and the behavior of a robot
at the moment of contact is determined by software that is increasingly
machine-generated. In a disembodied pipeline an error yields a wrong number a
human can catch downstream; in an embodied surgical system the same error
class yields a wrong motion a human cannot catch in time. No Federal law
requires the verification of robot-patient interaction code to clear before
that code is generated or executed, because the order of operations is not
regulated \cite{kawchak2026vvuq01, kawchak2026national}.}

\uscprov{1}{\textbf{(c) The Financial Data Gap.} Federal law already makes the
adjacent money flows of a clinical investigation reportable: a clinical
investigator's financial arrangements must be certified or disclosed under
part 54 of title 21, Code of Federal Regulations, and an applicable
manufacturer's payments to covered recipients are reported and published
annually under the Open Payments system established by section 1128G of the
Social Security Act (42 U.S.C. 1320a-7h) \cite{cfr-part54, usc-open-payments,
cfr-403-subpart-i}. But no Federal law prices, records, or discloses what it
costs to verify robot-patient interaction code, so a sponsor, a reviewer, and
the Congress each lack the financial data needed to judge whether safety
assurance is being bought, skipped, or padded. This Act closes that gap with
the financial-data record in section 515D(k) and the transparency duties of
section 5.}

\uscprov{1}{\textbf{(d) The Cost Asymmetry.} An embodied failure is priced in
lives, litigation, and program termination; an automated ex ante gate is
priced in compute, storage, and a bounded amount of human review. The
asymmetry is the financial case for the sequencing rule: the gate's unit cost
is small, fixed, and measurable, while the failure it prevents is large,
unbounded, and borne by patients first. Table 1 attaches the illustrative
financial footprint to each of the four works so the asymmetry can be read
from the record itself.}

{\footnotesize
\begin{xltabular}{\textwidth}{@{}L{1.5cm} Y L{4.7cm} R{1.7cm} L{1.7cm}@{}}
\hline\noalign{\vskip 0.04cm}
\textbf{Work} & \textbf{What it is} & \textbf{Headline assurance result} &
\textbf{Illustr.\ cost (\$)} & \textbf{End-goal DOI}\\
\hline\noalign{\vskip 0.04cm}
\endfirsthead
\hline\noalign{\vskip 0.04cm}
\textbf{Work} & \textbf{What it is} & \textbf{Headline assurance result} &
\textbf{Illustr.\ cost (\$)} & \textbf{End-goal DOI}\\
\hline\noalign{\vskip 0.04cm}
\endhead
\hline\noalign{\vskip 0.04cm}
\endlastfoot
VVUQ-01 & Method paper plus working pipeline and an execution record & 51 of
51 tests; 1 ACCEPT, 5 BLOCK, 1 ESCALATE gate surface & 12,600 &
\href{https://doi.org/10.5281/zenodo.20372501}{20372501}\\
\hline\noalign{\vskip 0.04cm}
VVUQ-02 & Ten-gate humanoid codebase, autonomous
execution record & 172 of 172 tests; 32 of 32 sweep iterations; composite mean
93.56 & 21,400 &
\href{https://doi.org/10.5281/zenodo.20421754}{20421754}\\
\hline\noalign{\vskip 0.04cm}
VVUQ-03 & Physical AI Oncology Trial Bill from the
recorded evidence & 15 sections; 21 body-width tables; 60 references & 6,100 &
\href{https://doi.org/10.5281/zenodo.20454870}{20454870}\\
\hline\noalign{\vskip 0.04cm}
VVUQ-04 & Bill as amendment to the Federal Food, Drug, and Cosmetic
Act & New section 360e-5 plus ten conforming amendments & 8,300 &
\href{https://doi.org/10.5281/zenodo.20485580}{20485580}\\
\hline\noalign{\vskip 0.04cm}
Total & Four-work evidence-to-law record & \raisebox{0.1ex}\textasciitilde 11 days of
auto. real-time work & 48,400 & -\\
\hline
\end{xltabular}}
\vspace{-0.75cm}
\figcaption{Table 1. The four works at a glance, with the financial footprint:
each work's role, headline\\assurance result, illustrative all-in build cost
(compute, storage, and records, from the seed\\20260525), and DOI. The
illustrative column sums to the 48,400 dollar total carried in Figure 2.}

\uscprov{1}{\textbf{(e) Feasibility Is Demonstrated, and Affordable.}
Automated verification of robot-patient interaction code is feasible today: a
pipeline held the verification fraction at 1.0, and a ten-gate assurance
suite over control software for a surgical humanoid cleared reproducibly from
a fixed seed, each gate bound to a published external standard
\cite{kawchak2026vvuq01, kawchak2026vvuq02, asmevv40, nasastd7009}. The same
record shows what the assurance cost: Figure 2 reads the conventional
program-year ledger beside the four-work record, and the appendices price the
unit economics gate by gate.}

\begin{asciifig}
Conventional path (per program-year)        Four-work autonomous record
+-----------------------------------+       +-----------------------------------+
| Manual code review   $1,800,000/y |       | VVUQ-01  pipeline      ~5 days    |
| Serial V&V cycles      $640,000/y |  vs   | VVUQ-02  10 gates      ~2 days    |
| Re-review after each              |       | VVUQ-03  bill          ~1 day     |
|   change             $310,000/y   |       | VVUQ-04  amendment     ~2 days    |
| Counsel + drafting   $250,000+    |       | compute + records   $48,400 total |
+-----------------------------------+       +-----------------------------------+
    ~$3.0 million per program-year              ~$0.05 million, fixed seed
    (illustrative, fully loaded)                (illustrative, seed 20260525)
\end{asciifig}
\vspace{-0.75cm}
\figcaption{Figure 2. Conventional versus autonomous cost ledger (ASCII). The
illustrative\\fully loaded conventional ledger beside the four-work record
whose per-work figures\\appear in Table 1; Appendix B reconciles the two
columns at program-year terms.}

\uscprov{1}{\textbf{(f) The Closest Existing Analogues.} The predetermined
change control plan authority in section 515C of the Federal Food, Drug, and
Cosmetic Act (21 U.S.C. 360e-4) and the Food and Drug Administration final
guidance on such plans for artificial-intelligence-enabled device software
functions are the closest existing analogues to validating a change before
deployment, but each remains voluntary, submission-based, and neither
automated nor required in advance of generation; the Quality Management System
Regulation and part 11 of title 21, Code of Federal Regulations, supply
quality and record-integrity scaffolding but no rule that verification
precedes generation and no cost record for the verification itself
\cite{fda-pccp-2024, fr-qmsr-2024, cfr-part11}. Human-over-AI review is
already widely adopted in State law, and the Medicare Advantage
artificial-intelligence guardrail is the strongest Federal precedent for
mandatory human oversight \cite{ca-sb1120, il-hb1806, tx-hb149, md-hb820,
cms-4201-ma}.}

\uscprov{1}{\textbf{(g) A Recognized Standards Basis.} Recognized consensus
standards, including ASME V\&V 40-2018, the Food and Drug Administration
guidance assessing the credibility of computational modeling, NASA-STD-7009A,
IEEE Std 7009-2024, IEC 80601-2-77, ISO 14971, and ISO/TS 15066, supply the
credibility framework, dispersion bounds, fail-safe requirements, and
biomechanical contact limits to which an automated gate may be bound, and the
graduated autonomy taxonomy maintained by the American Medical Association
provides the spine for matching verification rigor, and therefore
verification cost, to the degree of autonomy of a device \cite{asmevv40,
fda-credibility-2023, nasastd7009, ieee7009, iec8060127, iso14971, isots15066,
cpt-appendix-s}.}

\uscprov{1}{\textbf{(h) Budget Discipline.} The requirement rides the
financial machinery that already exists. Implementation is discretionary
spending within the Food and Drug Administration's appropriated device
program, supplemented by the medical-device user-fee structure of sections
737 and 738 (21 U.S.C. 379i, 379j) as it stands; this Act creates no
entitlement, no new tax, and no new statutory fee, authorizes appropriations
in stated amounts in section 5, and is therefore consistent with the
Statutory Pay-As-You-Go Act of 2010 and the House CutGo rule, as the cost
appendix demonstrates \cite{usc-379i, usc-379j, usc-paygo-2010,
house-rule-xxi, gao-budget-glossary}.}
